% Draft of the supplement subsection for study IH.  To use it once experiments/results/ih_summary.json exists:
%   cp docs/paper/ih_section.draft.tex docs/paper/ih_section.tex  (then edit the "Reading" paragraph to what the numbers show).  The include is added to supplement.tex by \IfFileExists.
\subsection{Study IH: A Destination Rule for the Request Parser}\label{sec:moved-ih}
Study IG found requests for a place the mission does not know (``to the community hall'', ``to the city hall'', ``to shelter 1'', ``to the shelter on the north side'') read as requests for shelter B. Study IH (registered and stamped before any inference on its texts) tests a fail-closed rule written from the principle ``a named place other than B'': a text that names a place noun from a fixed list, an oriented or numbered other shelter, or an English place is refused before the model is asked (default off: the submitted parser is unchanged; a text that names shelter B is exempt from the place-noun list). The set is \IHTexts{} fresh texts, 40 of them wrong destinations, 20 of those naming places that are \emph{not} on the rule's list on purpose.
\emph{Claims.} IH1: at most 5 unsafe outcomes with the rule, \IHUnsafe{} (Wilson upper bound \IHUnsafeUpper; submitted parser: \IHUnsafeV): \IHOneHolds. IH2: at least 80\% of the legitimate texts accepted with the right supply, \IHAccepted{} (\IHAcceptedRate; submitted parser: \IHAcceptedV): \IHTwoHolds. IH3: fewer unsafe outcomes than the submitted parser: \IHThreeHolds. Wrong-destination texts accepted: \IHWrongDest{} with the rule (\IHWrongDestV{} without); by kind (unsafe of the kind): listed places \IHListedBad, unlisted places \IHUnlistedBad, oriented or numbered shelters \IHOrientedBad. Predictions: at most 8 wrong-destination texts accepted \IHPredOneHolds; at most 4 of the 8 place-context texts refused (\IHContextBad) \IHPredTwoHolds. Post hoc, the IG answers with the rule on top: \IHIgUnsafe{} unsafe, \IHIgAccepted{} legitimate texts accepted.
\emph{Reading.} IH1 fails, as the registration predicted for the places the list does not contain. The rule does what it was written for: all ten wrong destinations on its list are refused (the submitted parser accepted five of them), and the unsafe outcomes fall from \IHUnsafeV{} to \IHUnsafe{}. What remains are nine places that the list does not contain (a mountain hut, a parking lot, a bus stop, a stadium, a golf course, a farmers' cooperative office, a place under a bridge, a municipal pool, a shopping street) and ``the second shelter'' written with a kanji ordinal that the patterns miss. The model's own destination question does not help: for these ten texts its probability of ``other'' (0.02--0.35) lies inside the range of the accepted legitimate texts (up to 0.44), so no threshold on it separates them. The price of the rule is small and in the safe direction: of the legitimate texts it falsely refuses \IHContextBad{} of the place-context texts (a place word that is only context or a source, with no shelter named) and one negation, each costing the operator a question. A list is therefore not a solution; the operator's read-back of the parsed destination before departure (deployment note) stays the control, and the rule is kept as an optional, default-off layer. A generative extraction of the destination phrase, or a positive allow-list of destinations, would be a new parser and needs a new fresh set. Limits: 133 texts written by the same assistant as the rule; the unlisted places are a sample of one author's imagination; Japanese and English only.
